\documentclass[12pt,largemargins]{homework}

\begin{document}
    
\question#{10}


ماشین حالت متناهی (\lr{FSM}) زیر از یک تراشه استخراج شده است و قصد داریم زبان (پروتکل ارتباطی) پذیرفته‌شده توسط آن را بیابیم. با استفاده از الگوریتم حذف وضعیت‌ها (\lr{State Elimination}) و رعایت دقیق ترتیب حذف از چپ به راست ($c_i, b_i$)، عبارت منظم (\lr{RE}) معادل این ماشین را به‌دست آورید.

\vspace{1cm}

\begin{center}
\begin{latin}
\begin{tikzpicture}[shorten >=1pt, on grid, auto,
                    every state/.style={draw, minimum size=0.9cm, text=black}]
                    
  % Center Node
  \node[state] (a) at (0,0) {$a$};
  
  % Diagonal initial arrow
  \draw[->, >=stealth, thick] (-1.2, 1.2) -- (a.135);

  % Branch 1 (Left)
  \node[state, accepting] (b1) at (-2.5, 0) {$b_1$};
  \node[state]            (c1) at (-2.5, 2.5) {$c_1$};

  % Branch 2 (Top)
  \node[state, accepting] (b2) at (0, 2.5) {$b_2$};
  \node[state]            (c2) at (2.5, 2.5) {$c_2$};

  % Branch 3 (Right)
  \node[state, accepting] (b3) at (2.5, 0) {$b_3$};
  \node[state]            (c3) at (2.5, -2.5) {$c_3$};

  % Branch 4 (Bottom)
  \node[state, accepting] (b4) at (0, -2.5) {$b_4$};
  \node[state]            (c4) at (-2.5, -2.5) {$c_4$};

  % Transitions
  \path[->]
    % Center to inner branches
    (a) edge node [below] {0} (b1)
    (a) edge node [right] {1} (b2)
    (a) edge node [below] {0} (b3)
    (a) edge node [right] {1} (b4)

    % Branch 1 (Left) b1 <-> c1
    (b1.70)  edge node [right] {0} (c1.290)
    (c1.250) edge node [left]  {1} (b1.110)
    (b1)     edge [loop left]  node {1} ()
    (c1)     edge [loop above] node {1} ()

    % Branch 2 (Top) b2 <-> c2
    (b2.340) edge node [below] {0} (c2.200)
    (c2.160) edge node [above] {1} (b2.20)
    (b2)     edge [loop above] node {1} ()
    (c2)     edge [loop right] node {0} ()

    % Branch 3 (Right) b3 <-> c3
    (b3.290) edge node [right] {0} (c3.70)
    (c3.110) edge node [left]  {1} (b3.250)
    (b3)     edge [loop right] node {1} ()
    (c3)     edge [loop below] node {1} ()

    % Branch 4 (Bottom) b4 <-> c4
    (b4.200) edge node [below] {0} (c4.340)
    (c4.20)  edge node [above] {1} (b4.160)
    (b4)     edge [loop below] node {1} ()
    (c4)     edge [loop left]  node {0} ();

\end{tikzpicture}
\end{latin}
\end{center}


\ifcompileanswers\solution

\begin{center}
\begin{latin}
\begin{tikzpicture}[shorten >=1pt, auto,
                    every state/.style={draw, minimum size=0.8cm, text=black}]
                    
  % State S (Gray)
  \node[state, fill=gray!30] (S) at (-5.5, -1.5) {$S$};

  % Center State a
  \node[state] (a) at (0, 0) {$a$};

  % The b_i States
  \node[state] (b1) at (-3, 0) {$b_1$};
  \node[state] (b2) at (0, 3) {$b_2$};
  \node[state] (b3) at (3, 0) {$b_3$};
  \node[state] (b4) at (0, -3) {$b_4$};

  % The c_i States
  \node[state] (c1) at (-3, 3) {$c_1$};
  \node[state] (c2) at (3, 3) {$c_2$};
  \node[state] (c3) at (3, -3) {$c_3$};
  \node[state] (c4) at (-3, -3) {$c_4$};

  % Final State F (Gray)
  \node[state, accepting, fill=gray!30] (F) at (11, 0) {$F$};

  % Transitions
  \path[->]
    % S to a
    (S) edge [out=0, in=225] node [above left] {$\varepsilon$} (a)

    % a to b_i
    (a) edge node [below] {$0$} (b1)
    (a) edge node [left] {$1$} (b2)
    (a) edge node [above] {$0$} (b3)
    (a) edge node [right] {$1$} (b4)

    % b1 and c1
    (b1) edge [bend right=15] node [right] {$0$} (c1)
    (c1) edge [bend right=15] node [left] {$1$} (b1)
    (b1) edge [loop left] node {$1$} ()
    (c1) edge [loop above] node {$1$} ()

    % b2 and c2
    (b2) edge [bend right=15] node [below] {$0$} (c2)
    (c2) edge [bend right=15] node [above] {$1$} (b2)
    (b2) edge [loop above] node {$1$} ()
    (c2) edge [loop above] node {$0$} ()

    % b3 and c3
    (b3) edge [bend right=15] node [left] {$0$} (c3)
    (c3) edge [bend right=15] node [right] {$1$} (b3)
    (b3) edge [loop above] node {$1$} ()
    (c3) edge [loop below] node {$1$} ()

    % b4 and c4
    (b4) edge [bend right=15] node [above] {$0$} (c4)
    (c4) edge [bend right=15] node [below] {$1$} (b4)
    (b4) edge [loop below] node {$1$} ()
    (c4) edge [loop left] node {$0$} ()

    % b_i to F (Epsilon transitions)
    (b3) edge node [above] {$\varepsilon$} (F)
    (b1) edge [out=70, in=110, distance=7cm] node [above] {$\varepsilon$} (F)
    (b2) edge [out=45, in=120, distance=4.5cm] node [below] {$\varepsilon$} (F)
    (b4) edge [out=-45, in=-90, distance=6cm] node [below] {$\varepsilon$} (F);

\end{tikzpicture}
\end{latin}
\end{center}


\begin{center}
\begin{latin}
\begin{tikzpicture}[shorten >=1pt, auto,
                    every state/.style={draw, minimum size=0.8cm, text=black}]
                    
  % State S 
  \node[state, fill=gray!30] (S) at (-5.5, -1.5) {$S$};

  % Center State a
  \node[state] (a) at (0, 0) {$a$};

  % The b_i States
  \node[state] (b1) at (-3, 0) {$b_1$};
  \node[state] (b2) at (0, 3) {$b_2$};
  \node[state] (b3) at (3, 0) {$b_3$};
  \node[state] (b4) at (0, -3) {$b_4$};

  % Final State F 
  \node[state, accepting , fill=gray!30] (F) at (11, 0) {$F$};

  % Transitions
  \path[->]
    % S to a
    (S) edge [out=0, in=225] node [above left] {$\varepsilon$} (a)

    % a to b_i
    (a) edge node [below] {$0$} (b1)
    (a) edge node [left] {$1$} (b2)
    (a) edge node [above] {$0$} (b3)
    (a) edge node [right] {$1$} (b4)

    % Loops on b_i
    (b1) edge [loop left] node {$1 + 01^+$} ()
    (b2) edge [loop above] node {$1 + 0^+1$} ()
    (b3) edge [loop above] node {$1 + 01^+$} ()
    (b4) edge [loop below] node {$1 + 0^+1$} ()

    % b_i to F (Epsilon transitions)
    (b3) edge node [above] {$\varepsilon$} (F)
    (b1) edge [out=90, in=100, distance=7cm] node [above] {$\varepsilon$} (F)
    (b2) edge [out=0, in=135, distance=3cm] node [above] {$\varepsilon$} (F)
    (b4) edge [out=15, in=200, distance=3cm] node [below] {$\varepsilon$} (F);

\end{tikzpicture}
\end{latin}
\end{center}

\begin{center}
\begin{latin}
\begin{tikzpicture}[shorten >=1pt, auto,
                    every state/.style={draw, minimum size=0.8cm, text=black}]
                    
  % ==========================================
  % Part 1 (Top Graph)
  % ==========================================
  % States
  \node[state, fill=gray!30] (S1) at (-5.5, -1.5) {$S$};
  \node[state] (a1) at (0, 0) {$a$};
  \node[state] (b3) at (2.5, 0) {$b_3$};
  \node[state] (b4) at (0, -2.5) {$b_4$};
  \node[state, accepting, fill=gray!30] (F1) at (10, 0) {$F$};

  % Transitions
  \path[->]
    (S1) edge [out=0, in=225] node [above left] {$\varepsilon$} (a1)
    (a1) edge node [above] {$0$} (b3)
    (a1) edge node [right] {$1$} (b4)
    (b3) edge [loop above] node {$1 + 01^+$} ()
    (b4) edge [loop below] node {$1 + 0^+1$} ()
    (b3) edge node [above] {$\varepsilon$} (F1)
    (b4) edge node [below] {$\varepsilon$} (F1);

  % ==========================================
  % Part 2 (Middle Graph)
  % ==========================================
  % States
  \node[state, fill=gray!30] (S2) at (-5.5, -6.5) {$S$};
  \node[state] (a2) at (0, -5) {$a$};
  \node[state, accepting, fill=gray!30] (F2) at (10, -5) {$F$};

  % Transitions
  \path[->]
    (S2) edge [out=0, in=225] node [above left] {$\varepsilon$} (a2)
    (a2) edge node [above] {$0(1 + 01^+)^* + 1(1 + 0^+1)^*$} (F2);

  % ==========================================
  % Part 3 (Bottom Graph)
  % ==========================================
  % States
  \node[state, fill=gray!30] (S3) at (-4, -9) {$S$};
  \node[state, accepting, fill=gray!30] (F3) at (8.5, -9) {$F$};

  % Transitions
  \path[->]
    (S3) edge node [above] {$RE = 0(1 + 01^+)^* + 1(1 + 0^+1)^*$} (F3);

\end{tikzpicture}
\end{latin}
\end{center}

\fi\end{document}